%%%%%%%%%%%%%%%%%%%%%%%%%%%%%%%%%
%
%       EXERCÍCIO
%
%%%%%%%%%%%%%%%%%%%%%%%%%%%%%%%%%

\definevalorquestao

\begin{exercicioBanco}[\valorquestao]
Cientistas brasileiros verificaram que uma determinada colônia de bactérias triplica a cada meia hora. Uma amostra de \(\n{10000}\) bactérias por mililitro foi colocada em um tubo de ensaio e, após um tempo \(x\), verificou-se que o total era de \(2,43 \times 10^6\) bactérias por mililitro. O valor de \(x\) é:

% Define as alternativas
\newcommand{\alternativas}{%
    \begin{center}
        \begin{tabularx}{\textwidth}{XXX}
            (a) duas horas. &
            (b) duas horas e 30 minutos. &
            (c) 3 horas e trinta minutos. \\[5pt]
            (d) 48 horas. &
            (e) 264 horas.
        \end{tabularx}
    \end{center}
}

% Define a resposta correta
\newcommand{\resposta}{B}

% Lógica condicional para exibição
\ifdefstring{\atividade}{lista}{%
    \alternativas
    \vspace{0.5em}
    
    \noindent\textbf{Resposta:} letra \textbf{\resposta}.
}{%
    \ifdefstring{\modo}{objetivo}{%
        \alternativas
    }{%
        % modo = discursiva → não mostra alternativas
    }
}
\end{exercicioBanco}

